\documentclass[a4,11pt]{aleph-notas}
% Se puede ver la documentación aquí:
% https://github.com/alephsub0/LaTeX_aleph-notas

% -- Paquetes adicionales
\usepackage{enumitem}
\usepackage{aleph-comandos}
\usepackage{booktabs}
\input{codigo-python.tex}
 

% -- Datos
\institucion{Facultad de Ciencias Exactas, Naturales y Ambientales}
\carrera{Catálogo STEM}
\asignatura{Métodos Numéricos}
\tema{Resumen no. 6: Integración, errores y optimización}
\autor{Mario Cueva - Andrés Merino}
\fecha{Periodo 2026-1}

\logouno[0.18\textwidth]{Logos/logoPUCE_04_ac}
\definecolor{colortext}{HTML}{1957d1}
\definecolor{colordef}{HTML}{1957d1}
\fuente{montserrat}


% -- Comandos adicionales
\setlist[enumerate]{label=\roman*.}
\newcommand{\vect}[1]{\mathbf{#1}}


\begin{document}

\encabezado

%%%%%%%%%%%%%%%%%%%%%%%%%%%%%%%%%%%%%%
\section{Integración numérica: regla de Simpson}

\begin{defi}[Regla de Simpson $1/3$ simple]
    La \textbf{regla de Simpson $1/3$ simple} aproxima la integral de una función $f$ en $[a,b]$ usando una parábola que interpola los puntos
    \[
        (a,f(a)),\qquad \left(\frac{a+b}{2},f\left(\frac{a+b}{2}\right)\right),\qquad (b,f(b)).
    \]
    Si $m=\dfrac{a+b}{2}$, entonces
    \[
        \int_a^b f(x)\,dx
        \approx
        \frac{b-a}{6}\left[f(a)+4f(m)+f(b)\right].
    \]
    Esta fórmula es exacta para polinomios de grado menor o igual que $3$.
\end{defi}

\begin{ejem}
    Para aproximar
    \[
        \int_0^2 (x^2+1)\,dx
    \]
    mediante Simpson $1/3$ simple, se tiene $a=0$, $b=2$ y $m=1$. Entonces
    \[
        f(0)=1,\qquad f(1)=2,\qquad f(2)=5.
    \]
    Por tanto,
    \[
        \int_0^2 (x^2+1)\,dx
        \approx
        \frac{2}{6}\left[1+4(2)+5\right]
        =
        \frac{14}{3}.
    \]
    En este caso el resultado coincide con el valor exacto, pues $x^2+1$ es un polinomio de grado $2$.
\end{ejem}

\begin{defi}[Regla de Simpson $1/3$ compuesta]
    La \textbf{regla de Simpson $1/3$ compuesta} divide $[a,b]$ en $n$ subintervalos de igual longitud
    \[
        h=\frac{b-a}{n},
    \]
    donde $n$ debe ser par. Si $x_i=a+ih$ para $i=0,1,\ldots,n$, entonces
    \[
        \int_a^b f(x)\,dx
        \approx
        \frac{h}{3}
        \left[
            f(x_0)
            +4\sum_{\substack{i=1\\ i\text{ impar}}}^{n-1}f(x_i)
            +2\sum_{\substack{i=2\\ i\text{ par}}}^{n-2}f(x_i)
            +f(x_n)
        \right].
    \]
\end{defi}

\begin{advertencia}
    En la regla de Simpson $1/3$ compuesta, el número de subintervalos debe ser par. Si se elige un $n$ impar, los puntos no pueden agruparse en pares de subintervalos para formar las parábolas.
\end{advertencia}

\begin{pscodigo}
\begin{function}[H]
\KwData{Una función $f$, extremos $a$ y $b$, número par de subintervalos $n$}
\KwResult{Aproximación de $\displaystyle\int_a^b f(x)\,dx$}
\SetKwFunction{simp}{Simpson}
\Fn{\simp{$f,a,b,n$}}{
    \If{$n$ no es par}{
        \Return{``El número de subintervalos debe ser par''}
    }
    $h:=\dfrac{b-a}{n}$\;
    $S:=f(a)+f(b)$\;
    \For{$i:=1$ \KwTo $n-1$}{
        \eIf{$i$ es impar}{
            $S:=S+4\cdot f(a+i\cdot h)$\;
        }{
            $S:=S+2\cdot f(a+i\cdot h)$\;
        }
    }
    \Return{$\dfrac{h}{3}\cdot S$}
    }
\end{function}
\end{pscodigo}

\begin{ejem}
    Para aproximar
    \[
        \int_0^1 e^x\,dx
    \]
    con Simpson $1/3$ compuesta usando $n=4$, se tiene $h=0.25$ y los nodos $0,0.25,0.5,0.75,1$:
    \[
        \begin{array}{c|ccccc}
            x_i & 0 & 0.25 & 0.5 & 0.75 & 1\\
            \hline
            e^{x_i} & 1.0000 & 1.2840 & 1.6487 & 2.1170 & 2.7183
        \end{array}
    \]
    Entonces
    \[
        \int_0^1 e^x\,dx
        \approx
        \frac{0.25}{3}
        \left[
            1.0000+4(1.2840+2.1170)+2(1.6487)+2.7183
        \right]
        \approx
        1.7183.
    \]
    El valor exacto es $e-1\approx 1.7183$, por lo que el error es muy pequeño.
\end{ejem}

\begin{defi}[Regla de Simpson $3/8$ simple]
    La \textbf{regla de Simpson $3/8$ simple} aproxima la integral usando un polinomio cúbico que interpola cuatro puntos igualmente espaciados. Si
    \[
        h=\frac{b-a}{3},
    \]
    entonces
    \[
        \int_a^b f(x)\,dx
        \approx
        \frac{3h}{8}\left[f(x_0)+3f(x_1)+3f(x_2)+f(x_3)\right],
    \]
    donde $x_0=a$, $x_1=a+h$, $x_2=a+2h$ y $x_3=b$.
\end{defi}

%%%%%%%%%%%%%%%%%%%%%%%%%%%%%%%%%%%%%%
\section{Errores de estimación de diferenciación e integración}

\begin{defi}[Error de truncamiento]
    El \textbf{error de truncamiento} aparece cuando se reemplaza un proceso exacto, como una derivada o una integral, por una fórmula aproximada que usa un número finito de términos o evaluaciones.
\end{defi}

En diferenciación numérica, las fórmulas de diferencias finitas se obtienen a partir de desarrollos de Taylor. Por ello, su error depende de potencias de $h$ y de derivadas de orden superior.

\begin{center}
\small
\begin{tabular}{p{0.26\textwidth}p{0.42\textwidth}p{0.18\textwidth}}
\toprule
Fórmula & Aproximación & Error \\
\midrule
Diferencia hacia adelante & $\dfrac{f(x+h)-f(x)}{h}$ & $O(h)$ \\[4mm]
Diferencia hacia atrás & $\dfrac{f(x)-f(x-h)}{h}$ & $O(h)$ \\[4mm]
Diferencia centrada & $\dfrac{f(x+h)-f(x-h)}{2h}$ & $O(h^2)$ \\[4mm]
Segunda derivada centrada & $\dfrac{f(x+h)-2f(x)+f(x-h)}{h^2}$ & $O(h^2)$ \\[2mm]
\bottomrule
\end{tabular}
\end{center}

\begin{defi}[Error de integración numérica]
    En integración numérica, el error depende de la longitud de los subintervalos y de la suavidad de la función. Para reglas compuestas, al disminuir $h$ se espera reducir el error, aunque con más evaluaciones de la función.
\end{defi}

\begin{center}
\small
\begin{tabular}{p{0.22\textwidth}p{0.46\textwidth}p{0.18\textwidth}}
\toprule
Método & Fórmula compuesta & Error global \\
\midrule
Trapecio & $\dfrac{h}{2}\left[f(x_0)+2\displaystyle\sum_{i=1}^{n-1}f(x_i)+f(x_n)\right]$ & $O(h^2)$ \\[6mm]
Simpson $1/3$ & $\dfrac{h}{3}\left[f(x_0)+4\displaystyle\sum_{i\text{ impar}}f(x_i)+2\displaystyle\sum_{i\text{ par}}f(x_i)+f(x_n)\right]$ & $O(h^4)$ \\[7mm]
\bottomrule
\end{tabular}
\end{center}

\begin{advertencia}
    Un método de orden más alto suele ser más preciso cuando la función es suficientemente suave. Sin embargo, si los datos tienen ruido o si $h$ es demasiado pequeño, los errores de redondeo pueden afectar el resultado.
\end{advertencia}


\begin{ejem}
    Para $f(x)=e^x$ en $x=1$ con $h=0.1$, la diferencia centrada da
    \[
        f'(1)\approx \frac{e^{1.1}-e^{0.9}}{0.2}
        =
        \frac{3.0042-2.4596}{0.2}
        =
        2.723.
    \]
    Como el valor exacto es $e\approx 2.718$, el error absoluto es aproximadamente
    \[
        |2.718-2.723|=0.005.
    \]
\end{ejem}

\begin{ejem}
    Para aproximar $\displaystyle\int_0^1 e^x\,dx$ con $n=4$, la regla del trapecio compuesta produce aproximadamente $1.7272$, mientras que Simpson $1/3$ compuesta produce aproximadamente $1.7183$. Como el valor exacto es
    \[
        e-1\approx 1.7183,
    \]
    Simpson obtiene un error menor en este caso.
\end{ejem}

\begin{defi}[Estimación por refinamiento]
    Si se calcula una aproximación con paso $h$ y otra con paso $h/2$, se puede estimar el error comparando ambos resultados:
    \[
        E_{\text{aprox}}=|I_{h/2}-I_h|.
    \]
    Esta idea permite evaluar la estabilidad de una aproximación aun cuando no se conoce el valor exacto.
\end{defi}


%%%%%%%%%%%%%%%%%%%%%%%%%%%%%%%%%%%%%%
\section{Optimización numérica}

\begin{defi}[Problema de optimización]
    Un problema de \textbf{optimización} busca encontrar un punto donde una función alcanza un valor mínimo o máximo. En una dimensión, se busca un valor $x^\ast$ que minimice o maximice $f(x)$. En varias variables, se busca un vector ${x}^\ast$ que optimice $f({x})$.
\end{defi}

\begin{defi}[Punto crítico]
    Si $f$ es derivable, un \textbf{punto crítico} en una dimensión es un valor $x^\ast$ tal que
    \[
        f'(x^\ast)=0.
    \]
    En varias variables, un punto crítico satisface
    \[
        \nabla f({x}^\ast)={0}.
    \]
\end{defi}

\subsection{Método de Newton para optimización en una dimensión}

Para encontrar puntos críticos de $f$, se puede aplicar Newton-Raphson a la ecuación
\[
    f'(x)=0.
\]
Así, si $f''(x_i)\neq 0$, la iteración queda
\[
    x_{i+1}=x_i-\frac{f'(x_i)}{f''(x_i)}.
\]

\begin{defi}[Newton para puntos críticos]
    El \textbf{método de Newton para optimización} aproxima un punto crítico de una función derivable dos veces mediante
    \[
        x_{i+1}=x_i-\frac{f'(x_i)}{f''(x_i)}.
    \]
    Si el método converge a $x^\ast$ y $f''(x^\ast)>0$, el punto es un mínimo local. Si $f''(x^\ast)<0$, es un máximo local.
\end{defi}

\begin{ejem}
    Consideremos
    \[
        f(x)=x^2-4x+5.
    \]
    Entonces
    \[
        f'(x)=2x-4,\qquad f''(x)=2.
    \]
    Con $x_0=0$,
    \[
        x_1=0-\frac{-4}{2}=2.
    \]
    Como $f'(2)=0$ y $f''(2)>0$, el punto $x=2$ es un mínimo local.
\end{ejem}


\subsection{Gradiente y dirección de máximo crecimiento}

\begin{defi}[Gradiente]
    Para una función $f(x_1,x_2,\ldots,x_n)$, el \textbf{gradiente} es el vector
    \[
        \nabla f({x})
        =
        \begin{pmatrix}
            \dfrac{\partial f}{\partial x_1}\\[2mm]
            \dfrac{\partial f}{\partial x_2}\\
            \vdots\\
            \dfrac{\partial f}{\partial x_n}
        \end{pmatrix}.
    \]
    El gradiente apunta en la dirección de máximo crecimiento local de la función.
\end{defi}

\begin{advertencia}
    Si se desea minimizar una función, la dirección natural de avance no es $\nabla f({x})$, sino $-\nabla f({x})$, porque esta dirección produce el mayor descenso local.
\end{advertencia}

\subsection{Funciones convexas y matriz hessiana}

\begin{defi}[Función convexa]
    Una función $f$ es \textbf{convexa} si, para cualesquiera puntos ${x}$ y ${y}$ de su dominio y para todo $t\in[0,1]$, se cumple
    \[
        f(t{x}+(1-t){y})
        \leq
        tf({x})+(1-t)f({y}).
    \]
    En una función convexa, todo mínimo local es también un mínimo global.
\end{defi}

\begin{defi}[Matriz hessiana]
    Si $f$ tiene derivadas parciales de segundo orden, su \textbf{matriz hessiana} es
    \[
        H_f({x})
        =
        \begin{pmatrix}
            \dfrac{\partial^2 f}{\partial x_1^2} &
            \dfrac{\partial^2 f}{\partial x_1\partial x_2} &
            \cdots &
            \dfrac{\partial^2 f}{\partial x_1\partial x_n}\\[3mm]
            \dfrac{\partial^2 f}{\partial x_2\partial x_1} &
            \dfrac{\partial^2 f}{\partial x_2^2} &
            \cdots &
            \dfrac{\partial^2 f}{\partial x_2\partial x_n}\\
            \vdots & \vdots & \ddots & \vdots\\
            \dfrac{\partial^2 f}{\partial x_n\partial x_1} &
            \dfrac{\partial^2 f}{\partial x_n\partial x_2} &
            \cdots &
            \dfrac{\partial^2 f}{\partial x_n^2}
        \end{pmatrix}.
    \]
    Si $H_f({x})$ es semidefinida positiva en una región convexa, entonces $f$ es convexa en esa región.
\end{defi}

\begin{ejem}
    Para
    \[
        f(x,y)=x^2+y^2,
    \]
    se tiene
    \[
        \nabla f(x,y)=
        \begin{pmatrix}
            2x\\
            2y
        \end{pmatrix},
        \qquad
        H_f(x,y)=
        \begin{pmatrix}
            2 & 0\\
            0 & 2
        \end{pmatrix}.
    \]
    La hessiana es definida positiva, por lo que la función es convexa y su mínimo global ocurre en $(0,0)$.
\end{ejem}

\subsection{Método de descenso del gradiente}

\begin{defi}[Descenso del gradiente]
    El \textbf{método de descenso del gradiente} es un método iterativo para minimizar una función diferenciable. A partir de una aproximación inicial ${x}^{(0)}$, actualiza
    \[
        {x}^{(k+1)}
        =
        {x}^{(k)}
        -
        \alpha \nabla f\bigl({x}^{(k)}\bigr),
    \]
    donde $\alpha>0$ es el \textbf{tamaño de paso} o \textbf{learning rate}.
\end{defi}

\begin{pscodigo}
\begin{function}[H]
\KwData{Función $f$, gradiente $\nabla f$, aproximación inicial ${x}^{(0)}$, learning rate $\alpha$, tolerancia $\varepsilon$ y máximo de iteraciones $N$}
\KwResult{Aproximación de un mínimo de $f$}
\SetKwFunction{grad}{DescensoGradiente}
\Fn{\grad{$f,\nabla f,{x}^{(0)},\alpha,\varepsilon,N$}}{
    \For{$k:=1$ \KwTo $N$}{
        ${x}^{(k)}:={x}^{(k-1)}-\alpha\nabla f({x}^{(k-1)})$\;
        \If{$\left\|{x}^{(k)}-{x}^{(k-1)}\right\|<\varepsilon$}{
            \Return{${x}^{(k)}$}
        }
    }
    \Return{${x}^{(k)}$}
    }
\end{function}
\end{pscodigo}

\begin{advertencia}
    Si $\alpha$ es demasiado grande, el método puede oscilar o divergir. Si $\alpha$ es demasiado pequeño, el método puede converger muy lentamente. Elegir un buen learning rate es una parte central del método.
\end{advertencia}

\begin{ejem}
    Consideremos
    \[
        f(x,y)=x^2+y^2.
    \]
    Su gradiente es
    \[
        \nabla f(x,y)=
        \begin{pmatrix}
            2x\\
            2y
        \end{pmatrix}.
    \]
    Con punto inicial ${x}^{(0)}=(2,1)$ y $\alpha=0.25$, la primera iteración es
    \[
        {x}^{(1)}
        =
        \begin{pmatrix}
            2\\
            1
        \end{pmatrix}
        -
        0.25
        \begin{pmatrix}
            4\\
            2
        \end{pmatrix}
        =
        \begin{pmatrix}
            1\\
            0.5
        \end{pmatrix}.
    \]
    La siguiente iteración produce
    \[
        {x}^{(2)}
        =
        \begin{pmatrix}
            0.5\\
            0.25
        \end{pmatrix}.
    \]
    Las aproximaciones se acercan al mínimo global $(0,0)$.
\end{ejem}


\subsection{Learning rate dinámico}

\begin{defi}[Learning rate dinámico]
    Un \textbf{learning rate dinámico} cambia el tamaño de paso durante las iteraciones. En lugar de usar un valor fijo $\alpha$, se usa una sucesión
    \[
        \alpha_0,\alpha_1,\alpha_2,\ldots
    \]
    y la actualización toma la forma
    \[
        {x}^{(k+1)}
        =
        {x}^{(k)}
        -
        \alpha_k \nabla f\bigl({x}^{(k)}\bigr).
    \]
    La idea usual es empezar con pasos relativamente grandes y reducirlos gradualmente para estabilizar la convergencia.
\end{defi}

Algunas estrategias comunes son:
\begin{enumerate}
    \item \textbf{Decaimiento constante:} reducir $\alpha$ cada cierto número de iteraciones.
    \item \textbf{Decaimiento proporcional:} usar una regla como $\alpha_k=\dfrac{\alpha_0}{1+\lambda k}$.
    \item \textbf{Búsqueda de línea:} elegir $\alpha_k$ probando qué tamaño de paso reduce mejor la función objetivo.
\end{enumerate}

\begin{advertencia}
    Un learning rate dinámico no garantiza convergencia por sí solo. Su utilidad depende de la función, de la escala de las variables y de la forma en que se ajusta $\alpha_k$ durante el proceso.
\end{advertencia}

\begin{defi}[Métodos adaptativos]
    Los \textbf{métodos adaptativos} ajustan el tamaño de paso de forma diferente para cada variable, usando información acumulada de gradientes anteriores. Entre los más conocidos están AdaGrad, RMSProp y Adam.
\end{defi}

\begin{center}
\small
\begin{tabular}{p{0.22\textwidth}p{0.62\textwidth}}
\toprule
Método & Idea principal \\
\midrule
Momentum & Acumula parte de la dirección anterior para suavizar el movimiento \\
AdaGrad & Reduce más el paso en variables con gradientes acumulados grandes \\
RMSProp & Usa un promedio móvil de gradientes cuadrados para ajustar el paso \\
Adam & Combina ideas de momentum y RMSProp \\
\bottomrule
\end{tabular}
\end{center}

\begin{ejem}
    Si se usa un learning rate dinámico de la forma
    \[
        \alpha_k=\frac{0.1}{1+0.05k},
    \]
    entonces
    \[
        \alpha_0=0.1,\qquad
        \alpha_{10}=\frac{0.1}{1.5}\approx 0.0667,\qquad
        \alpha_{100}=\frac{0.1}{6}\approx 0.0167.
    \]
    Al inicio el método permite pasos más grandes; luego los pasos se reducen para mejorar la estabilidad cerca del mínimo.
\end{ejem}






\end{document}
